% --- 1. INFORMAZIONI GENERALI ---
\section{Informazioni Generali}
\begin{itemize}
	\item \textbf{Azienda Coinvolta:} Zucchetti
	\item \textbf{Mezzo di Comunicazione:} Email
	\item \textbf{Data:} 2026-04-15
\end{itemize}

Il presente verbale documenta uno scambio di email tra il gruppo Synergo Unit
e il referente dell'azienda Zucchetti, avvenuto nell'ambito della stesura
dell'Analisi dei Requisiti per il progetto \textit{Second Brain}.
L'obiettivo della comunicazione è stato ottenere chiarimenti su tre aspetti
emersi durante i lavori interni di approfondimento: le funzionalità
dell'editor Markdown, le lingue da supportare per la traduzione tramite
LLM\textsubscript{G} e il profilo dell'utenza target del prodotto.

\vspace{0.5cm}
\textbf{Referenti Azienda:}
\begin{table}[ht]
	\centering
	\renewcommand{\arraystretch}{1.2}
	\begin{tabularx}{\textwidth}{X l}
		\toprule
		\textbf{Nome e Cognome} & \textbf{Ruolo} \\
		\midrule
		Gregorio Piccoli & Referente Principale \\
		\bottomrule
	\end{tabularx}
\end{table}

% --- 2. CHIARIMENTI TECNICI RICHIESTI ---
\section{Chiarimenti Richiesti}
Il gruppo ha richiesto chiarimenti sui seguenti aspetti:

\begin{enumerate}
	\item Funzionalità dell'editor e gestione del Markdown.
	\begin{itemize}
		\item Completezza del set di strumenti di formattazione individuato
		      (grassetto, corsivo, sottolineato, barrato, titoli, link
		      ipertestuali, elenchi puntati e numerati, tabelle, citazioni).
		\item Eventuale distinzione di priorità tra funzionalità obbligatorie
		      e facoltative.
		\item Inclusione o esclusione della gestione delle immagini come
		      requisito del progetto.
	\end{itemize}
	\item Integrazione con i modelli linguistici per la funzionalità di
	      traduzione del testo.
	\begin{itemize}
		\item Lingue da supportare.
		\item Numero di idiomi previsti dai requisiti.
	\end{itemize}
	\item Utenza target del prodotto.
	\begin{itemize}
		\item Presenza di un profilo utente specifico o adozione di un
		      profilo generico.
	\end{itemize}
\end{enumerate}

% --- 3. RESOCONTO ---
\section{Resoconto}
Dalla risposta del referente aziendale sono emersi i seguenti punti.

\subsection{Funzionalità Dell'Editor E Gestione Del Markdown}
\begin{itemize}
	\item Il set di funzionalità di formattazione individuato dal gruppo
	      è ritenuto completo e adeguato; tutte le funzionalità elencate
	      devono essere incluse nel prodotto.
	\item La gestione delle immagini non è considerata un requisito del
	      progetto e deve essere esclusa.
\end{itemize}

\subsection{Lingue Supportate Per La Traduzione}
\begin{itemize}
	\item L'obiettivo non è realizzare un sistema di traduzione professionale,
	      ma dimostrare la capacità di eseguire il compito tramite LLM.
	\item Le lingue da supportare sono: Inglese, Francese e Spagnolo; il
	      Tedesco può essere incluso facoltativamente.
	\item Sono escluse le lingue con alfabeti non latini (es.\ Cinese,
	      Russo, Arabo), in quanto la loro gestione introdurrebbe problematiche
	      di presentazione del testo (es.\ scrittura da destra a sinistra)
	      che esulano dagli obiettivi del progetto.
\end{itemize}

\subsection{Utenza Target}
\begin{itemize}
	\item Il profilo utente di riferimento è un utente giovane, tipicamente
	      uno studente, che utilizza l'applicazione per prendere appunti e
	      desidera essere supportato nella loro elaborazione tramite LLM.
\end{itemize}

% --- 4. DECISIONI FORMALI ---
\section{Decisioni Formali}
Nella seguente tabella sono tracciate le decisioni adottate a seguito
dei chiarimenti ricevuti.

\renewcommand{\arraystretch}{1.3}
\vspace{-0.3cm}
\noindent
\begin{longtable}{c p{12cm}}
	\toprule
	\textbf{Codice} & \textbf{Decisione} \\
	\midrule
	\endfirsthead

	\toprule
	\textbf{Codice} & \textbf{Decisione} \\
	\midrule
	\endhead

	\midrule
	\multicolumn{2}{r}{\textit{Continua nella pagina successiva}} \\
	\midrule
	\endfoot

	\bottomrule
	\endlastfoot

	\textbf{VE-5.1}
		& Includere nell'editor tutte le funzionalità di formattazione
		  individuate (grassetto, corsivo, sottolineato, barrato, titoli,
		  link ipertestuali, elenchi puntati e numerati, tabelle, citazioni).
		  La gestione delle immagini è esclusa dai requisiti. \\[4pt]

	\textbf{VE-5.2}
		& La funzionalità di traduzione tramite LLM supporterà le lingue
		  Inglese, Francese e Spagnolo come requisito obbligatorio; il
		  Tedesco potrà essere incluso come estensione facoltativa. Sono
		  escluse le lingue con alfabeti non latini. \\[4pt]

	\textbf{VE-5.3}
		& Il profilo utente di riferimento per la definizione dei casi
		  d'uso\textsubscript{G} è un utente giovane (es.\ studente) che
		  utilizza l'applicazione per prendere appunti e desidera supporto
		  nella loro elaborazione tramite LLM. \\

\end{longtable}

\newpage
% --- 5. FIRME ---
\section*{Approvazione e Firme}
Con le seguenti firme, i rappresentanti approvano formalmente il contenuto
del presente verbale e le decisioni adottate.

\vspace{2.5cm}

\noindent
\begin{minipage}[c]{0.45\textwidth}
	\begin{center}
		\rule{6cm}{0.4pt} \\
		\vspace{0.2cm}
		\textbf{Per il Gruppo Synergo Unit} \\
		Sofia De Blasi \\
		\textit{Responsabile di Progetto}
	\end{center}
\end{minipage}
\hfill
\begin{minipage}[c]{0.45\textwidth}
	\begin{center}
		\rule{6cm}{0.4pt} \\
		\vspace{0.2cm}
		\textbf{Per l'Azienda Zucchetti} \\
		Gregorio Piccoli \\
		\textit{Referente Aziendale}
	\end{center}
\end{minipage}

\vspace{2cm}